\documentclass[a4paper]{article}

\usepackage{graphicx}
\usepackage{amsmath,amssymb}
\usepackage{float}
\usepackage{natbib}

\title{Proportional integral derivative controller}
\date{}
\author{Mohammad Rahmani}

\begin{document}
   \maketitle

   A proportional--integral--derivative (PID) controller is a feedback controller that combines the present control error, the accumulated past error, and the rate of change of the error to generate a corrective control input \cite{astrom-2008-feedback-systems-an-introduction-for-scientists-and-engineers}.

   \section{Formulation}
      Let the reference value be $r(t)$ and the measured system output be $y(t)$. The control error is
      \[
         e(t)=r(t)-y(t).
      \]

      The continuous-time PID control law is
      \[
         u(t)
         =
         K_P e(t)
         +
         K_I \int_{0}^{t} e(\tau)\,d\tau
         +
         K_D \frac{de(t)}{dt}.
      \]

      where $K_P$, $K_I$, and $K_D$ are the proportional, integral, and derivative gains.

   \subsection{Contribution of each term}
      \begin{itemize}
         \item \textbf{Proportional term:} $K_P e(t)$ reacts to the current error. A larger error produces a larger immediate corrective action.
         \item \textbf{Integral term:} $K_I \int_0^t e(\tau)\,d\tau$ accumulates persistent error and can remove steady-state offset, but excessive integral action can produce overshoot and integral windup.
         \item \textbf{Derivative term:} $K_D \frac{de(t)}{dt}$ reacts to how quickly the error is changing and can provide damping, but it is sensitive to measurement noise.
      \end{itemize}
      These effects and the practical trade-offs among the three terms are standard properties of PID control \cite{astrom-2008-feedback-systems-an-introduction-for-scientists-and-engineers}.

   \subsection{Vector form}
      For a multi-dimensional controlled variable,
      \[
         \mathbf{e}(t)
         =
         \mathbf{r}(t)-\mathbf{y}(t),
      \]
      and
      \[
         \mathbf{u}(t)
         =
         K_P\mathbf{e}(t)
         +
         K_I\int_0^t\mathbf{e}(\tau)\,d\tau
         +
         K_D\dot{\mathbf{e}}(t),
      \]
      where the gains may be scalars or matrices, depending on whether the controlled dimensions are coupled.

   \subsection{Discrete-time form}
      For sampled control with sampling interval $\Delta t$,
      \[
         \mathbf{I}_k
         =
         \mathbf{I}_{k-1}
         +
         \mathbf{e}_k\Delta t,
         \qquad
         \mathbf{D}_k
         =
         \frac{\mathbf{e}_k-\mathbf{e}_{k-1}}{\Delta t},
      \]
      and the PID command is
      \[
         \mathbf{u}_k
         =
         K_P\mathbf{e}_k
         +
         K_I\mathbf{I}_k
         +
         K_D\mathbf{D}_k.
      \]

   \subsection{Velocity-control loop}
      In a cascaded position/velocity controller, the outer position loop can produce a velocity setpoint $\mathbf{v}^{\mathrm{sp}}_k$. The inner PID loop then uses
      \[
         \mathbf{e}_{v,k}
         =
         \mathbf{v}^{\mathrm{sp}}_k-\mathbf{v}_k
      \]
      and generates the low-level command
      \[
         \mathbf{c}_k
         =
         K_{\mathrm{P},v}\mathbf{e}_{v,k}
         +
         K_{\mathrm{I},v}
         \sum_{j=0}^{k}\mathbf{e}_{v,j}\Delta t
         +
         K_{\mathrm{D},v}
         \frac{\mathbf{e}_{v,k}-\mathbf{e}_{v,k-1}}{\Delta t}.
      \]
      Here $\mathbf{c}_k$ denotes the low-level command produced by the velocity controller. This notation is useful when a higher-level controller already uses $\mathbf{u}_k$ for a supervisory control input.

   \subsection{Transfer-function form}
      The transfer function of an ideal PID controller is
      \[
         C(s)
         =
         K_P
         +
         \frac{K_I}{s}
         +
         K_D s.
      \]

      For a plant with transfer function $G(s)$ and unity negative feedback, the closed-loop transfer function is
      \[
         \frac{Y(s)}{R(s)}
         =
         \frac{C(s)G(s)}
              {1+C(s)G(s)}.
      \]

   \section{Practical considerations}
      \begin{itemize}
         \item Larger $K_P$ increases the immediate corrective response, but excessive proportional gain can increase oscillation or instability.
         \item Larger $K_I$ increases correction of persistent bias and steady-state error, but excessive integral action can cause windup when the actuator saturates.
         \item Larger $K_D$ increases damping against rapid error changes, but derivative action amplifies high-frequency measurement noise and is therefore commonly filtered in practical implementations.
      \end{itemize}

      PID gains must therefore be chosen with respect to the plant dynamics, sampling interval, actuator limits, and measurement noise \cite{astrom-2008-feedback-systems-an-introduction-for-scientists-and-engineers}.

   \bibliographystyle{plainnat}
   \bibliography{/home/donkarlo/Dropbox/repo/nd_shared_research/bibliography/bibliography.bib}
\end{document}
